\section{مقدمه}

باتری‌های لیتیوم-یون به دلیل هزینه نسبتاً پایین و رو به کاهش، چگالی انرژی بالا و طول عمر مناسب، در طیف گسترده‌ای از کاربردها مورد استفاده قرار می‌گیرند. با این حال، این باتری‌ها در اثر استفاده مداوم و تکرار چرخه‌های شارژ و دشارژ\LTRfootnote{Charge/Discharge Cycles}، به‌تدریج دچار افت ظرفیت و کاهش عملکرد می‌شوند. این فرایند که با عنوان تخریب باتری\LTRfootnote{Battery Degradation} شناخته می‌شود، می‌تواند بر قابلیت اطمینان، ایمنی و طول عمر سیستم تأثیر بگذارد. همچنین، همان‌طور که در بسیاری از سیستم‌های شیمیایی، مکانیکی و الکترونیکی صادق است، ارزیابی باتری‌هایی با طول عمر بالا به آزمون‌های طولانی‌مدت نیاز دارد؛ به‌طوری‌که این فرایند ممکن است چندین ماه یا حتی چند سال به طول بینجامد.

پیش‌بینی عمر مفید باقی‌مانده باتری‌های لیتیوم-یون با استفاده از مدل‌های مبتنی بر فیزیک، حتی برای باتری‌های یک تولیدکننده و در شرایط کاری یکسان، چالش‌برانگیز است؛ زیرا افت ظرفیت باتری‌ها در اثر چرخه‌های شارژ و دشارژ، رفتاری غیرخطی دارد. در این میان، روش‌های مبتنی بر یادگیری عمیق که قابلیت استخراج خودکار الگوها و ویژگی‌های پیچیده از داده‌های اندازه‌گیری‌شده را دارند، می‌توانند رویکرد مناسبی برای تخمین عمر مفید باقی‌مانده باتری‌ها فراهم آورند.

در این گزارش، روش ارائه‌شده در مثال \lr{Battery Cycle Life Prediction Using Deep Learning} از مستندات \lr{MathWorks} \cite{mathworks} بررسی می‌شود. در این مثال، از داده‌های اندازه‌گیری‌شده مربوط به چرخه‌های شارژ و دشارژ باتری‌های لیتیوم-یون برای آموزش یک شبکه عصبی کانولوشنی\LTRfootnote{Convolutional Neural Network (CNN)} دوبعدی استفاده شده است. هدف از این روش، تخمین تعداد چرخه‌های باقی‌مانده عمر باتری با استفاده از شبکه آموزش‌دیده و ارزیابی عملکرد مدل است.
